\section{Proofs and auxiliary lemmas}

This appendix records the auxiliary probability result used by the large-sample analysis: the normal approximation for bounded summands indexed by a sparse dependency graph.

\subsection{Dependency-graph normal approximation}

The linearization summands in \cref{obj:def:first-order-linearization} are dependent only when their outcome neighborhoods overlap. The next lemma states the bounded-degree central-limit result used to convert that structure into the asymptotic normality asserted in \cref{obj:thm:hetero-clt}; it is the bounded-summand, bounded-dependency-degree specialization of the dependency-graph normal approximation of \citet{BaldiRinott1989}.

\begin{lemmav}{lem:bounded-degree-dependency-clt}[Bounded-Degree Dependency CLT]
Let \(\{X_{n,i}: i\in\iota_n\}\) be real-valued random variables, where \(|\iota_n|=n\), under the usual measurability and integrability conditions. Suppose:
\begin{itemize}
\item \textbf{(Centered summands.)} For every \(n\) and \(i\in\iota_n\), \(\mathbb E[X_{n,i}]=0\).
\item \textbf{(Dependency degree.)} The variables admit a dependency graph whose neighborhood of every \(i\in\iota_n\) has cardinality at most \(D\), for a fixed \(D\in\mathbb N\).
\item \textbf{(Uniform bound.)} There is an \(M\ge 0\) such that \(|X_{n,i}|\le M\) for every \(n\), \(i\in\iota_n\), and outcome.
\item \textbf{(Variance growth.)} Writing \(S_n=\sum_{i\in\iota_n}X_{n,i}\), suppose \(\mathbb E[S_n^2]=v_n\) for every \(n\), and, for some \(c>0\), \(v_n\ge cn\) for all sufficiently large \(n\).
\end{itemize}
Then, for every \(s\in\mathbb R\),
\[
\lim_{n\to\infty}\Pr\!\left(\frac{S_n}{\sqrt{v_n}}\le s\right)
=
\Pr\!\left(Z\le s\right),
\qquad Z\sim\mathcal N(0,1).
\]
\end{lemmav}

The imported hypotheses map directly to the paper's primitives. Bounded outcomes, bounded outcome degree, and the positivity floor give a uniform bound on each centered linearization summand. The outcome-overlap graph is a dependency graph because disjoint intervention neighborhoods depend on disjoint independent assignments, and \cref{obj:ass:bounded-overlap-dependency} bounds its closed-neighborhood size. \Cref{obj:ass:variance-nondegenerate}, together with the paper's normalization, supplies variance growth of order \(n\). These are precisely the boundedness, local-dependence, and nondegeneracy inputs used in the displayed design-based specialization of \citet{BaldiRinott1989}.

\subsection{Convex-program provenance}

The KKT characterization used in \cref{obj:thm:convex-design} is the standard convex-program optimality system for a differentiable objective with affine budget and box constraints \citep[Chs.~4--5]{BoydVandenberghe2004}. The positivity box keeps the reciprocal objective differentiable; the theorem supplies convexity and existence; the budget equality and coordinate bounds are affine. Its displayed multipliers state dual feasibility, stationarity, and complementary slackness coordinate by coordinate. The verified theorem directly includes feasible endpoint budgets and the corresponding optimality conditions.

\paragraph{Verification note.}
The theorem and lemma statements, together with their proofs for the finite-design and independent Bernoulli-randomization components, are machine-checked in Lean~4. The formal development takes as hypotheses the experiment-specific bipartite graph and neighborhood structure, the fixed potential-outcome schedule and interference restriction, feasibility and positivity conditions, and the stated asymptotic regularity conditions on outcomes, degrees, overlap dependence, indexing, and variance nondegeneracy. Classical mathematical and probability-theory inputs used by the underlying libraries form the assumed formal substrate.
